\documentclass{article}
\usepackage[utf8]{inputenc}
\usepackage[T1]{fontenc}
\usepackage{geometry}
\geometry{
    a4paper,
    left=2.5cm,
    right=2.5cm,
    top=2.5cm,
    bottom=2.5cm
}
\usepackage{graphicx}
\usepackage{float}
\usepackage{hyperref}
\usepackage[english]{babel}
\usepackage{amsmath}
\usepackage{amssymb}
\usepackage{array}
\usepackage{booktabs}
\usepackage{caption}
\usepackage{subcaption}
\usepackage{enumitem}
\usepackage{fancyhdr}
\usepackage{titlesec}
\usepackage{url}
\usepackage{textcomp}
\usepackage{xcolor}
\usepackage{listings}
\usepackage{multirow}
\usepackage{longtable}
\usepackage{tabularx}
\usepackage{times}

% Set up page style
\pagestyle{fancy}
\fancyhf{}
\rfoot{\thepage}
\renewcommand{\headrulewidth}{0pt}

% Title formatting
\titleformat{\chapter}[display]{\normalfont\bfseries}{\chaptertitle\thechapter}{0pt}{\Huge}
\titlespacing*{\chapter}{0pt}{50pt}{40pt}

% Document info
\title{Template Document}
\author{Academic Support Agent}
\date{\today}

\begin{document}
\maketitle

\begin{titlepage}
\begin{figure}[!htb]
\includegraphics[width=15.92cm, height=4.03cm]{logo_PPGEE.png}
\label{logo}
\end{figure}
\begin{center}
\par
\vspace{100pt}
\textbf{\Large Nome do Aluno}\\
\vspace{102pt}
\textbf{\Large T tulo}\\
\end{center}

\par
\vfill
\begin{center}
{\singlespacing{\normalsize Jo o Pessoa - PB}\\
{\normalsize Mar o de \the\year}}
\end{center}
\end{titlepage}

\thispagestyle{empty}

\newpage
\begin{figure}[!htb]
\includegraphics[width=15.92cm, height=4.03cm]{logo_PPGEE.png}
\label{logo2}
\end{figure}
\begin{center}
\vspace{30pt}
\textbf{\LARGE Nome do Aluno}

\par
\vspace{50pt}
\textbf{\Large T tulo}
\end{center}

\par
\vspace{30pt}
\hspace*{175pt}\parbox{7.6cm}{\singlespacing {\normalsize Exame de Qualifica o de Mestrado submetida ao Programa de P s-Gradua o em
Engenharia El trica do Instituto Federal de Educa o, Ci ncia e Tecnologia da Para ba, como requisito necess rio   obten o do grau de Mestre em Ci ncias no Dom nio da Engenharia El trica.}}

\par
\vspace{40pt}
\hspace*{125pt}\parbox{10.0cm}{{\normalsize  rea de Concentra o: Eletromagnetismo Avan ado}}

\par
\vspace{30pt}
\hspace*{125pt}\parbox{10.0cm}{{\normalsize Orientador: Nome do Orientador, Prof. Dr.}\\
{\normalsize Co-Orientador: Nome do Coorientador, Prof. Dr.}}

\par
\vfill
\begin{center}
{\singlespacing{\normalsize Jo o Pessoa - PB, Mar o de \the\year}}\\
{{\normalsize \copyright Nome do Aluno -  aluno.estudante@academico.ifpb.edu.com}}
\end{center}

% \thispagestyle{empty}
%  \newpage
%  \noindent{\textbf{\large Agradecimentos}}
%  \vspace{1pt}

% Agradecimentos...

% \thispagestyle{empty}
%  \newpage
%     \vspace*{\fill}
% 	\begin{flushright}
% 		\textit{``Nada lhe pertence mais que seus sonhos.''\\
% 		(Friedrich Nietzsche)}
% 	\end{flushright}

\thispagestyle{empty}
\newpage
\begin{center}
\textbf{\normalsize Resumo}
\end{center}
\vspace{1pt}
   
\noindent Escreva o seu resumo aqui.
\\  \\
Palavras-Chave: Modelo, latex, overleaf.

\thispagestyle{empty}

\newpage
\begin{center}
\textbf{\normalsize Abstract}
\end{center}
\vspace{1pt}

\noindent The english version of your "Resumo".
\\  \\
Keywords: Template, latex, overleaf.

\thispagestyle{empty}
\newpage

\newpage

\newpage

\newpage

% Content from cap1 would be included here
 
%% Content from cap2 would be included here
 
%% Content from cap3 would be included here
 
%% Content from cap4 would be included here
 
%% Content from cap5 would be included here
 
%% Content from cap6 would be included here


\newpage
\centering

 



% \newpage
% \section{Ap ndices}

% \subsection{xx}


% \begin{figure}[H]
% %left bottom right top
% %\includegraphics[trim=-10mm 55mm 00mm 15mm, scale=0.5]{gr fico.pdf}

% \begin{center}
% Figura 3: gr fico.
% \end{center}
% \end{figure}

% \begin{figure}[H]
% %left bottom right top
% %\includegraphics[trim=-10mm 55mm 00mm 15mm, scale=0.5]{gr fico.pdf}

% \caption{gr fico.}

% \end{figure}

% \newpage{}
% \subsection{Algoritmo}

% \subsubsection{yy}
% \begin{verbatim}
%   %Simula o;

%   for i = 1:n;
%     M(:,i) = rand(n,1);
%   end

% \end{verbatim}
\end{document}